\documentclass[UTF8,a4paper,12pt]{ctexart}
\usepackage[margin=2.5cm]{geometry}
\usepackage{fontspec,booktabs,tabularx,array,enumitem}
\usepackage[hidelinks]{hyperref}
\setmainfont{Times New Roman}
\setCJKmainfont{SimSun}[BoldFont=SimHei]
\setCJKsansfont{SimHei}
\setlength{\parindent}{2em}
\linespread{1.22}
\pagestyle{plain}
\ctexset{section={format=\large\bfseries,number=\chinese{section},aftername=、,beforeskip=14pt,afterskip=8pt}}
\setlist[enumerate]{leftmargin=2em,itemsep=5pt,topsep=4pt,parsep=0pt}
\newcommand{\field}[2]{\par\noindent\textbf{#1：}#2\par}
\newcommand{\example}[2]{\vspace{6pt}\noindent\textbf{#1、#2}\par}
\begin{document}
\begin{center}\LARGE\bfseries AI工具使用详情\end{center}
\vspace{3pt}

\section{AI工具名称和版本}
\begin{center}
\small 表1\quad AI工具名称、型号及使用范围\par\vspace{5pt}
\renewcommand{\arraystretch}{1.35}
\begin{tabularx}{\linewidth}{@{}p{2.9cm}p{3.5cm}X@{}}
\toprule
工具名称 & 版本或型号 & 具体使用目的和环节\\
\midrule
OpenAI Codex & GPT-6（按本项目会话标识；精确部署子版本未留存） & 题意梳理、数据整理、候选方案讨论、程序编写与执行、验证分析、图表绘制、论文起草与修订、参考文献核查、提交材料整理。\\
\bottomrule
\end{tabularx}
\end{center}

本报告依据2026年9月10日至13日本项目工作聊天及保存的使用记录编写，记录对象为《微网与外部电网电力调控策略》。Codex中的辅助审查任务属于同一工具的使用，不另列为一种AI。Python、SciPy、scikit-learn、Matplotlib和LaTeX为AI调用的计算或排版软件，不作为独立生成式AI工具列报。

\section{具体使用目的与环节}
\begin{enumerate}[label=\arabic*、]
\item \textbf{题意与假设梳理。}读取题面和附件，解释四问的信息条件、购电及储能关系、费用规则和提交内容；提出有歧义的口径，由队员讨论并明确采用方式。
\item \textbf{数据预处理与探索性分析。}生成附件读取和格式整理代码，核对时间标签、单位和原始数值；统计日内、星期及季节变化，生成数据分析图。分析发现用于支持后续预测方案讨论。
\item \textbf{候选讨论、程序实现与计算。}协助提出并比较候选方案，编写预测、线性规划、运行回放及费用计算程序，并按队员指令运行。四问的模型路线由队员结合解释与比较结果逐步选择。
\item \textbf{模型深化与结果核对。}协助设计并执行获准的参数比较、连续偏差检验和信息价值对照；编写核验程序，回算供需平衡、储电量、购电量与费用，保留无改善和费用增加的结果。
\item \textbf{论文和图表。}依据已保存研究材料起草章节、摘要、模型评价及通俗讲解，并根据队员反馈修改。使用Matplotlib、TikZ生成图表和框架图，调整重叠文字、图例及排版。
\item \textbf{引用与使用记录。}检查参考文献题录及公开来源，整理本报告和论文中的AI使用声明。AI参与范围包含方案建议、代码和初稿生成，不仅限于语法检查或语言润色。
\end{enumerate}

\newpage
\section{主要提示方式与使用过程说明}
\field{主要提示方式}{队员通过中文多轮对话提出任务，上传题面、附件、组内参考材料和论文截图；AI读取项目文件，生成或修改代码、运行程序并解释结果。涉及模型选择、运行口径和后续工作范围时，队员给出明确意见。}
以下交互根据本项目真实聊天节选或摘要整理，回答为内容概述，并非逐字完整对话。

\example{1}{供电关系与时间口径}
\field{队员提示}{“光伏可以直接给小区负载供电，不必先经过电池”；指出部分附件时间整体偏移10分钟，要求读取和填写结果时统一处理，并明确充电、放电效率各为90\%。}
\field{AI处理}{解释光伏直供、外网充电、逐段供需与储电递推的关系；将时间标签映射和两侧效率写入代码与记录。}
\field{队伍处理}{队员对物理关系、时间口径和效率作出明确选择，并多次追问其含义；不把这些运行假设冒充题面逐字规定。}

\example{2}{问题一候选比较与费用差异}
\field{队员提示}{“实际比较优先围绕3、4、10展开”；随后要求核查35126.85元与35126.95元的差异，并将附件第一行暂按0:00--0:10另算一次。}
\field{AI处理}{比较峰谷策略、线性规划和光伏预留策略；通过辅助审查与隔离时间对照，解释不同时间映射对应的0.10元差异。}
\field{队伍处理}{队员选择线性规划，并明确保留原时间理解及35126.85元版本；对照计算不用于替换正式方案。程序复算属于AI执行的核对。}

\example{3}{问题二分别预测与权重解释}
\field{队员提示}{“光伏和负载是要分开建模的”；要求记录权重分配方法，并指出仅运行150天不足，需要补齐全年，再补充验证。}
\field{AI处理}{分别建立供需预测并进行候选比较，解释近期误差准备、日期衰减及五倍紧急购电的关系；完成全年计算、参数比较与连续偏差分析。}
\field{队伍处理}{队员提出分别预测、全年运行及补充验证要求，审核论文框架后要求起草；AI被要求说明参数依据，不能只给费用排名。}

\example{4}{问题三的调整费用与预报价值}
\field{队员提示}{提出原计划100度电减至80度电时，退还减少部分货款并收违约费的理解；要求相对零点计划处理重叠调整，解释小时预报如何对应10分钟结果。}
\field{AI处理}{梳理原计划、调整和紧急购电的结算关系，生成预报补间、融合、动态准备与滚动规划代码；区分日内反馈和后发预报带来的费用变化；按后续要求加入近期同星期负载加权和日内偏差校正，再完成三组信息条件的同口径比较。}
\field{队伍处理}{队员确认口径、选择方案2作为基础，并要求解释权重和分位参数；写作中另行要求明确回答日内预报价值问题。}

\newpage
\noindent\textbf{三、主要提示方式与使用过程说明（续）}

\example{5}{问题四两条路线的选择}
\field{队员提示}{明确“问题4的电价不能提前知道”，随后要求“问题四第二问选方案3，第三问选方案2，开始设计深化与验证”。}
\field{AI处理}{使用历史可得信息预测电价；比较零点购电与日内调整两条分支，分别实现提升树和岭回归结合线性规划的方案，并补充特征、参数及风险分析；问题三更新后，重新运行问题4-3并核对全年结果。}
\field{队伍处理}{队员明确未来电价的信息限制和两条所选路线，批准深化后再要求论文框架及正文；并要求解释岭回归的作用与选择理由。}

\example{6}{论文结构与图表修正}
\field{队员提示}{要求论文按模型假设、符号说明、模型建立与求解组织；指出文字像研究报告、模型选择解释不足，以及三张图存在文字与曲线重叠。}
\field{AI处理}{重组并起草正文，根据队员反馈补充解释；从原Matplotlib代码重绘图表，把图例移出数据区、下移子图，保留原曲线和数值。}
\field{队伍处理}{队员通过阅读草稿和图片提出具体修改意见，明确删除内容、调整标题及布局。队员还要求保留必要解释，压缩重复文字与图表、统一求解步骤及约束标题。当前版本的最终采用以队伍提交文件为准。}

\section{对AI输出的采纳、人工修改与核验情况}

\textbf{采纳与人工决策。}聊天记录可确认的队员参与包括：提出并确认供电、效率、时间映射、跨日储电及费用口径；选择问题一线性规划，要求问题二分别预测并补齐全年；选择问题三融合预报路线，以及问题四两条不同价格预测路线。队员也对初稿结构、解释深度和图表布局进行了持续审阅与修改指示。上述行为是具体的决策和内容审阅，不等同于已对每个数值独立验算。

\textbf{AI执行的核对。}AI编写或运行了原附件与整理数据的逐值比对、信息时点检查、逐段供需和储电约束核查、费用回算、结果表回读、参数对照及PDF页面检查。9月13日按队员要求整理支撑材料，导出五个完整结果工作簿，并独立复算整理后的最终程序。问题一费用差异另使用辅助AI任务进行复核。此类程序检查及另一AI的复核均属于工具核对，本报告不将其记作队员人工验算。

\textbf{人工核验的记录边界。}现有记录能够证明队员对上述口径、模型选择、文本和图片提出过具体意见，但没有完整记录所有核心公式、代码和结果已逐项完成人工终审。提交前仍需由队员逐项审查与核实AI参与的核心内容。本报告不作“所有AI输出均已由队员逐项核验”的无依据声明。

\textbf{报告覆盖范围。}本报告汇总的是当前项目工作聊天中可追溯的Codex使用。组员在其他聊天中形成的参考材料及外部提供的立体示意图，制作工具和版本未在本记录中完整留存，不能据此断言未使用其他AI；全队提交版本需合并其实际使用信息。DeepSeek的实际使用型号、日期、用途和交互内容仍需团队补充；当前官网型号不作为实际使用的证明。

\end{document}
